% ── Section 6: Cross-case synthesis ─────────────────────────────
\section{Cross-Case Synthesis}
\label{sec:synthesis}

Chlorophyll and
dissolved oxygen respond to overlapping but different physical and biological processes, have
different distributions, and fail in different regimes. The transferable object is therefore the
benchmark workflow itself: define a daily target, construct artificial gaps, enforce leakage control and
compare methods against withheld truth. Under that workflow, the two targets show a consistent methodological
hierarchy but different scientific failure modes.

\subsection{What transfers across targets}
\label{sec:synthesis_transfers}

The first transferable result is that interpolation is a strong floor. In both case studies, simple
target continuity explains a large fraction of the recoverable signal, especially at short gaps. Methods that add
external predictors or model complexity must therefore be judged asking whether they improve on the information already carried by the local target trajectory.

The second transferable result is that external tabular prediction alone is not enough. For both
chlorophyll and oxygen, the external-only model is worse than interpolation. This means that, when
used as day-wise external predictors without target context, physical forcing data do not reliably replace the
local continuity of the in-situ record. In a nearshore bay, gridded products can carry useful physical
information while still being too coarse to reconstruct the sensor target by themselves.

The third transferable result is that pretrained sequence modelling adds skill beyond this floor.
TS-ICL is the strongest chlorophyll method on both benchmarks and this is operationally important because the model is used
zero-shot, which means its parameters are not fitted or fine-tuned on Tongoy data. 

\subsection{What changes between targets}
\label{sec:synthesis_changes}

The covariate evidence is target-specific. For chlorophyll, the clearest physical signal is
wind forcing, while the tested current inputs do not establish a useful
reconstruction signal. This is consistent with a biological target whose short events are linked to
upwelling and surface biomass but are also spatially patchy and difficult for coarse external products
to capture. For oxygen, current information behaves differently: it is informative in
isolation, but largely redundant when added to the broader physical covariate bundle. The same
predictor family can therefore matter differently depending on the target and on the other covariates
already present.

The main failure mode also changes. For chlorophyll, the hardest regime is high-\chl{} events, since every method has higher error on event gaps than on non-event gaps. For oxygen, the TS-ICL improvement
is concentrated in the middle of the empirical distribution, while the tails remain difficult. The same general lesson is that the method
that wins on average is not automatically reliable in the regimes that matter most scientifically.

This distinction is important for interpretation and therefore the benchmark needs both
pooled comparison and stratified diagnostics. Without the latter, the chlorophyll result
would overstate event skill and the oxygen one would have done the same with tail reliability.

\subsection{Implications for operational reconstruction}
\label{sec:synthesis_operational}

The practical implication of the results is the following. First, any candidate method should be benchmarked against
linear interpolation under artificial gaps. Second, external predictors should be tested as empirical
information channels and not assumed useful from physical plausibility alone. Third, a selected
reconstruction method should be accompanied by regime diagnostics.

This is where TS-ICL is operationally attractive. Once the target and covariate channels are
prepared, it does not require fitting a station-specific model for each sensor or rebuilding a
large engineered tabular design. Inference is fast, and the same framework can ingest target
history and covariates as time-indexed channels rather than forcing temporal awareness into a
large set of hand-built lags, rolling summaries, and gap-edge columns. The channel structure also
makes diagnostic tests straightforward: a covariate family can be added, removed, or replaced by
placebo versions to test whether it carries time-aligned information. Eventually, the reconstruction is not only a point estimate but a candidate trajectory
with an uncertainty band, which is a key factor in real gaps scenarios.